\chapter{Binomio y coeficientes multinómicos}
\label{md:binomio}

\section{Números combinatorios y triángulo de Pascal}

Para $0\le k\le n$,
\[
 \binom nk=\binom n{n-k},\qquad
 \binom n0=\binom nn=1.
\]
La simetría corresponde a elegir un subconjunto o su complementario.
La \textbf{identidad de Pascal}, para $1\le k\le n-1$, es
\[
 \binom nk=\binom{n-1}{k}+\binom{n-1}{k-1}.
\]
Se separan los subconjuntos de tamaño $k$ según contengan o no un
elemento fijado. Esta identidad permite construir el
\textbf{triángulo de Pascal}: los extremos de cada fila son $1$ y
cada entrada interior es la suma de las dos situadas encima.
La fila $n$ contiene $\binom n0,\ldots,\binom nn$.

\section{Teorema del binomio}

Para $n\in\N_0$ y $x,y\in\R$, y también para elementos de un cuerpo,
\[
 \boxed{(x+y)^n=\sum_{k=0}^n\binom nk x^{n-k}y^k.}
\]
El coeficiente cuenta en qué $k$ de los $n$ factores se elige $y$.
Otra demostración usa inducción y la identidad de Pascal.
Los exponentes de $x$ e $y$ suman $n$; el signo de una constante
sustituida en $y$ se conserva dentro de su potencia.

En particular, $(1+x)^n=\sum_{k=0}^n\binom nk x^k$.
Evaluar y, cuando procede, derivar esta identidad proporciona
\[
 \begin{aligned}
 \sum_{k=0}^n\binom nk&=2^n &&(n\ge0),\\
 \sum_{k=0}^n(-1)^k\binom nk&=0 &&(n\ge1),\\
 \sum_{k=1}^n k\binom nk&=n2^{n-1} &&(n\ge1).
 \end{aligned}
\]
La suma alternada dice que, para $n\ge1$, hay tantos subconjuntos
de cardinal par como impar.
Otra identidad útil, para $0\le k\le m$, es
\[
 \sum_{j=k}^m\binom jk=\binom{m+1}{k+1}.
\]
Se obtiene aplicando sucesivamente Pascal, o por inducción sobre $m$.

\begin{ejemplo}
El coeficiente de $x^2$ en $(x-2)^5$ es
$\binom52(-2)^3=-80$. El exponente de la constante es $5-2$,
lo que determina también el signo.
\end{ejemplo}

\section{Multinómicos y ordenaciones con elementos iguales}
\label{md:multinomial}

Si $n_1,\ldots,n_s\in\N_0$ y $n_1+\cdots+n_s=n$, se define
\[
 \binom{n}{n_1,\ldots,n_s}=\frac{n!}{n_1!\cdots n_s!}.
\]
Este \textbf{coeficiente multinómico} cuenta las ordenaciones lineales
de $n$ objetos de $s$ tipos, con $n_i$ objetos del tipo $i$,
indistinguibles dentro de cada tipo. Las posiciones siguen siendo distintas.
Se justifica eligiendo primero las $n_1$ posiciones del primer tipo,
luego las $n_2$ del segundo entre las restantes, y así sucesivamente.

La \textbf{fórmula multinómica}, o de Leibniz, generaliza el binomio:
para $n\ge0$ y elementos $x_1,\ldots,x_s$ de un cuerpo,
\[
 \boxed{(x_1+\cdots+x_s)^n=
 \sum_{\substack{n_1,\ldots,n_s\ge0\\n_1+\cdots+n_s=n}}
 \binom{n}{n_1,\ldots,n_s}x_1^{n_1}\cdots x_s^{n_s}.}
\]
Cada término distribuye las $n$ elecciones entre los $s$ sumandos;
el coeficiente cuenta las ordenaciones que originan el mismo monomio.
Para calcular un coeficiente se comprueba primero que los exponentes
son no negativos y suman $n$; un sumando ausente tiene exponente cero.
